% Part: proof-theory
% Chapter: normalization
% Section: reductions

\documentclass[../../../include/open-logic-section]{subfiles}

\begin{document}

\olfileid{pt}{nor}{red}

\olsection{تبدیل‌های کاهشی}

اگر برشی در !!a{proof} در \Log{N1i} طول~$1$ داشته باشد، قطعه‌ای شامل
تنها یک رخداد !!{formula} است. آن !!{formula} نتیجهٔ !!a{introduction}
یا نتیجهٔ~\FalseInt، و هم‌زمان مقدمهٔ عمدهٔ !!a{elimination} است. در
اثبات نرمال‌سازی، چنین برش‌هایی با جایگزینی زیر-!!{proof} منتهی به این
جفت !!{introduction}/!!{elimination} یا جفت
\FalseInt/حذف با !!{proof} تازه‌ای که این دو استنتاج در آن
حذف شده‌اند «کاهش» می‌یابند. این جایگزینی ممکن است برش‌های تازه‌ای
بسازد، اما برش‌های تازه رتبه‌ای کمتر از برش کاهش‌یافته دارند. هنگامی
که یک برش بیشینهٔ بالاترین و راست‌ترین را برمی‌گزینیم، !!{proof} تازه
یا طول برش کمتری دارد، چون یک برش بیشینهٔ طول~$1$ حذف شده، یا رتبهٔ
برشِ !!{proof} کاهش می‌یابد، اگر آن برش تنها برش بیشینه بوده باشد.

برای نمونه، اگر استنتاج‌های درگیر \Intro{\lif} و \Elim{\lif} باشند،
$!A\ident !B\lif !C$ و زیر-!!{proof}~$\delta_1$ای داریم که چنین پایان
می‌یابد:
\[
\AxiomC{$\Discharge{!B}{x}$}
\RightLabel{$\delta_2$}
\DeduceC{$!C$}
\DischargeRule{\Intro{\lif}}{x}
\UnaryInfC{$!B \lif !C$}
\AxiomC{}
\RightLabel{$\delta_3$}
\DeduceC{$!B$}
\RightLabel{\Elim{\lif}}
\BinaryInfC{$!C$}
\DisplayProof
\]
برای حذف این برش، زیر‌برهان~$\delta_1$ در~$\delta$ را با برهان زیر
جایگزین می‌کنیم:
\[
\AxiomC{}
\RightLabel{$\delta_3$}
\DeduceC{$!B$}
\RightLabel{$\Subst{\delta_2}{\delta_3}{!B^x}$}
\DeduceC{$!C$}
\DisplayProof
\]
منظور از $\Subst{\delta_2}{\delta_3}{!B^x}$، !!{proof}ی است که با
جایگزینی هر فرض~$!B$ دارای برچسب~$x$ در~$\delta_2$ با کل !!{proof}
~‌$\delta_3$ حاصل می‌شود. این !!{proof} دیگر فرض بازی به شکل~$!B$ با
برچسب~$x$ ندارد و هیچ !!{formula} دیگری نیز با همان برچسب ندارد. هر
نسخهٔ~$\delta_3$ را به‌گونه‌ای بی‌گرفتاری پیوند می‌زنیم: متغیرهای ویژه
و برچسب‌های ترخیص داخلی آن را نسبت به~$\delta_2$ و نسبت به دیگر نسخه‌ها
تازه می‌کنیم. از این رو !!{proof} تازه همان فرمول‌ها و برچسب‌های فرض
باز~$\delta_1$ را دارد، هرچند بعضی رخدادها ممکن است چند بار تکرار شوند؛
استنتاج‌های~$\delta_2$ همان فرض‌های پیشین را مرخص می‌کنند و هر نسخهٔ
استنتاج‌های مرخص‌کنندهٔ~$\delta_3$ فرض‌های متناظر همان نسخه را مرخص
می‌کند.

برش‌های سراسر در~$\delta_2$ و برش‌های سراسر بیرون از~$\delta_1$ در
~‌$\delta$ بی‌تغییر می‌مانند. برش‌های هر نسخهٔ تازه‌شدهٔ~$\delta_3$
نیز همان !!{formula}ها، رتبه و طول برش‌های متناظر در~$\delta_3$ را
دارند. انتخاب یک برش بیشینه که هم بالاترین و هم راست‌ترین است تضمین
می‌کند هیچ برش بیشینه‌ای در بخش‌های کپی‌شده وجود ندارد. پس با حذف برش
بیشینهٔ طول~$1$، یا طول برش کاهش می‌یابد، یا اگر این تنها برش بیشینه
بود، رتبهٔ برش کاهش می‌یابد.

با این همه دو رخداد ممکن است:
\begin{enumerate}
  \item با جایگزینی همهٔ فرض‌های~$!B^x$ با~$\delta_3$، همهٔ برش‌های
  ~‌$\delta_3$ را چند برابر کرده‌ایم. اگر برخی از آن‌ها بیشینه باشند،
  طول برشِ !!{proof} تازه ممکن است از طول برشِ !!{proof} پیشین بیشتر
  شود. باید جلوی این امر را بگیریم. از این رو تبدیل‌های کاهشی را فقط
  بر برش‌های بیشینه‌ای اعمال می‌کنیم که هم \emph{بالاترین} و هم
  \emph{راست‌ترین} هستند. اگر برشی بیشینه در~$\delta_3$ بود، در شاخه‌ای
  در سمت راست برش مورد کاهش قرار می‌گرفت؛ و اگر برشی بیشینه در بخش
  کپی‌شوندهٔ بالای آن بود، برش برگزیده بالاترین نبود. با این انتخاب یا
  طول برشِ !!{proof} کاهش می‌یابد یا حتی رتبهٔ برش کم می‌شود.

  \item اگر فرض~$!B^x$ در~$\delta_2$ مقدمهٔ عمدهٔ !!a{elimination}
  باشد و آخرین استنتاج~$\delta_3$، \Elim{\lor}، \Elim{\lexists}،
  !!a{introduction} یا~\FalseInt باشد، با پیوندزدن~$\delta_3$ به آن
  فرض در~$\delta_2$ برشی تازه ساخته‌ایم. چیزی که در~$\delta_2$ فقط یک
  فرض بود، اکنون برشی با طول~$1$ است---نتیجهٔ !!a{introduction} یا
  \FalseInt و مقدمهٔ عمدهٔ !!a{elimination}---یا آخرین رخداد
  !!{formula} یک قطعهٔ برشی، یعنی مقدمهٔ عمدهٔ !!a{elimination} و
  نتیجهٔ \Elim{\lor} یا \Elim{\lexists}. اگر $!B^x$ مقدمهٔ فرعی
  \Elim{\lor} یا \Elim{\lexists} بود، از پیش نخستین رخداد
  !!{formula} یک قطعه، و شاید قطعه‌ای برشی، بود. در !!{proof} تازه آن
  قطعه ممکن است بلندتر شود. اما !!{formula} سازندهٔ این برش تازه یا
  قطعهٔ برشی موجود~$!B$ است و
  $\depth{!B}<\depth{!B\lif !C}$. پس گرچه ممکن است برش افزوده یا طول
  برشی زیاد شده باشد، هیچ‌یک بیشینه نیست. بنابراین یا رتبهٔ برش را کم
  کرده‌ایم، یا اگر برش‌های هم‌رتبه باقی مانده‌اند، طول برش را کاهش
  داده‌ایم.
\end{enumerate}

فرایند کاهش برش با طول~$1$ را \emph{تبدیل کاهشی} یا \emph{تبدیل
انحراف} می‌نامند. الگوهای برخی جفت‌قاعده‌ها در
\olref{tab:reduction-conversions} آمده‌اند و بقیه به‌عنوان تمرین واگذار
می‌شوند.

\begin{table}
\begin{multline*}
\bottomAlignProof
\AxiomC{$\Discharge{!B}{x}$}
\RightLabel{$\delta_2$}
\shortDeduce
\DeduceC{$!C$}
\DischargeRule{\Intro{\lif}}{x}
\UnaryInfC{$!B \lif !C$}
\AxiomC{}
\RightLabel{$\delta_3$}
\shortDeduce
\DeduceC{$!B$}
\RightLabel{\Elim{\lif}}
\BinaryInfC{$!C$}
\DisplayProof
\quad \longrightarrow\quad
\shoveright{\bottomAlignProof
\AxiomC{}
\RightLabel{$\delta_3$}
\shortDeduce
\DeduceC{$!B$}
\RightLabel{$\Subst{\delta_2}{\delta_3}{!B^x}$}
\shortDeduce
\DeduceC{$!C$}
\DisplayProof}
\\[4ex]
\shoveleft{\bottomAlignProof
\AxiomC{}
\RightLabel{$\delta_2$}
\shortDeduce
\DeduceC{$!B$}
\AxiomC{}
\RightLabel{$\delta_3$}
\shortDeduce
\DeduceC{$!C$}
\RightLabel{\Intro{\land}}
\BinaryInfC{$!B \land !C$}
\RightLabel{\Elim{\land}[1]}
\UnaryInfC{$!B$}
\DisplayProof
\quad \longrightarrow\quad
\bottomAlignProof
\AxiomC{}
\RightLabel{$\delta_2$}
\shortDeduce
\DeduceC{$!B$}
\DisplayProof}\qquad
\shoveright{\bottomAlignProof
\AxiomC{}
\RightLabel{$\delta_2$}
\shortDeduce
\DeduceC{$!B$}
\AxiomC{}
\RightLabel{$\delta_3$}
\shortDeduce
\DeduceC{$!C$}
\RightLabel{\Intro{\land}}
\BinaryInfC{$!B \land !C$}
\RightLabel{\Elim{\land}[2]}
\UnaryInfC{$!C$}
\DisplayProof
\quad \longrightarrow\quad
\bottomAlignProof
\AxiomC{}
\RightLabel{$\delta_3$}
\shortDeduce
\DeduceC{$!C$}
\DisplayProof}
\\[4ex]
\bottomAlignProof
\AxiomC{}
\RightLabel{$\delta_2$}
\shortDeduce
\DeduceC{$!B$}
\RightLabel{\Intro{\lor}[1]}
\UnaryInfC{$!B \lor !C$}
\AxiomC{$\Discharge{!B}{x}$}
\RightLabel{$\delta_3$}
\shortDeduce
\DeduceC{$!D$}
\AxiomC{$\Discharge{!C}{y}$}
\RightLabel{$\delta_4$}
\shortDeduce
\DeduceC{$!D$}
\DischargeRule{\Elim{\lor}}{x\,y}
\TrinaryInfC{$!D$}
\DisplayProof
\quad \longrightarrow\quad
\shoveright{\bottomAlignProof
\AxiomC{}
\RightLabel{$\delta_2$}
\shortDeduce
\DeduceC{$!B$}
\RightLabel{$\Subst{\delta_3}{\delta_2}{!B^x}$}
\shortDeduce
\DeduceC{$!D$}
\DisplayProof}
\\[2ex]
\shoveleft{\bottomAlignProof
\AxiomC{}
\RightLabel{$\delta_2$}
\shortDeduce
\DeduceC{$!B(c)$}
\RightLabel{\Intro{\lforall}}
\UnaryInfC{$\lforall[x][!B(x)]$}
\RightLabel{\Elim{\lforall}}
\UnaryInfC{$!B(t)$}
\DisplayProof
\quad \longrightarrow\quad
\bottomAlignProof
\AxiomC{}
\RightLabel{$\Subst{\delta_2}{t}{c}$}
\shortDeduce
\DeduceC{$!B(t)$}
\DisplayProof}
\\[4ex]
\shoveright{\bottomAlignProof
\AxiomC{}
\RightLabel{$\delta_2$}
\shortDeduce
\DeduceC{$!B(t)$}
\RightLabel{\Intro{\lexists}}
\UnaryInfC{$\lexists[x][!B(x)]$}
\AxiomC{$\Discharge{!B(c)}{x}$}
\RightLabel{$\delta_3$}
\shortDeduce
\DeduceC{$!D$}
\DischargeRule{\Elim{\lexists}}{x}
\BinaryInfC{$!D$}
\DisplayProof
\quad \longrightarrow\quad
\bottomAlignProof
\AxiomC{}
\RightLabel{$\delta_2$}
\shortDeduce
\DeduceC{$!B(t)$}
\RightLabel{$\Subst{\Subst{\delta_3}{t}{c}}{\delta_2}{!B(t)^x}$}
\shortDeduce
\DeduceC{$!D$}
\DisplayProof}
\end{multline*}
\caption{چند تبدیل کاهشی برای \Log{N1i}.}
\ollabel{tab:reduction-conversions}
\end{table}

\begin{prob}
تبدیل کاهشی را برای \Intro{\lnot} که \Elim{\lnot} در پی آن می‌آید
بنویسید.
\end{prob}

یک حالت دیگر باقی می‌ماند. تعریف برش حالتی را نیز در بر می‌گیرد که طول
برش~$1$ است و~$!A$ هم نتیجهٔ~\FalseInt{} و هم مقدمهٔ عمدهٔ
!!a{elimination} است. این حالت مانند حالتی بررسی می‌شود که~$!A$ نتیجهٔ
!!a{introduction} است. برای نمونه، برهان زیر را
\[
\bottomAlignProof
\AxiomC{}
\RightLabel{$\delta_2$}
\DeduceC{$\lfalse$}
\RightLabel{\FalseInt}
\UnaryInfC{$!B \lif !C$}
\AxiomC{}
\RightLabel{$\delta_3$}
\DeduceC{$!B$}
\RightLabel{\Elim{\lif}}
\BinaryInfC{$!C$}
\DisplayProof
\quad\text{با}\quad
\bottomAlignProof
\AxiomC{}
\RightLabel{$\delta_2$}
\DeduceC{$\lfalse$}
\RightLabel{\FalseInt}
\UnaryInfC{$!C$}
\DisplayProof
\]
جایگزین کنید. در حالت‌های $!A\ident(!B\lor !C)$ یا
$!A\ident\lexists[x][!B(x)]$ باید کمی محتاط باشیم. برای نمونه، اگر
برهان
\[
\bottomAlignProof
\AxiomC{}
\RightLabel{$\delta_2$}
\DeduceC{$\lfalse$}
\RightLabel{\FalseInt}
\UnaryInfC{$!B \lor !C$}
\AxiomC{$\Discharge{!B}{x}$}
\RightLabel{$\delta_3$}
\DeduceC{$!D$}
\AxiomC{$\Discharge{!C}{y}$}
\RightLabel{$\delta_4$}
\DeduceC{$!D$}
\DischargeRule{\Elim{\lor}}{x\,y}
\TrinaryInfC{$!D$}
\DisplayProof
\quad\text{را با}\quad
\bottomAlignProof
\AxiomC{}
\RightLabel{$\delta_2$}
\DeduceC{$\lfalse$}
\RightLabel{\FalseInt}
\UnaryInfC{$!D$}
\DisplayProof
\]
جایگزین می‌کردیم، ممکن بود برش تازه‌ای با !!{formula} برش~$!D$ معرفی
کنیم و هیچ تضمینی نیست که
$\depth{!D}<\depth{!B\lor !C}$!{} پس باید همانند کاهش
\Intro{\lor}/\Elim{\lor} عمل کنیم و آن را با یکی از دو برهان زیر
جایگزین سازیم:
\[
\bottomAlignProof
\AxiomC{}
\RightLabel{$\delta_2$}
\DeduceC{$\lfalse$}
\RightLabel{\FalseInt}
\UnaryInfC{$!B$}
\RightLabel{$\delta_3$}
\DeduceC{$!D$}
\DisplayProof
\quad\text{یا}\quad
\bottomAlignProof
\AxiomC{}
\RightLabel{$\delta_2$}
\DeduceC{$\lfalse$}
\RightLabel{\FalseInt}
\UnaryInfC{$!C$}
\RightLabel{$\delta_4$}
\DeduceC{$!D$}
\DisplayProof
\]

\end{document}
